\documentclass[../../main.tex]{subfiles}% 注意这里的文件路径不能用 ./main.tex ,否则用latexmk编译子文件会报错
\graphicspath{{\subfix{./image/}}{\subfix{../../image/}}} % 指定图片引用路径,后续可以直接使用图片文件名
% 如果图片存放在上级目录,则使用 ../../image/ 作为备用路径

% 例如:
% \begin{figure}[H]
% \centering
% \includegraphics[scale=0.3]{图.png}
% \caption{}
% \label{figure:图}
% \end{figure}
% 注意:上述\label{}一定要放在\caption{}之后,否则引用图片序号会只会显示??.

\begin{document}

\section{模的直和}

\begin{definition}
设$M_1,M_2,\cdots,M_n$为$R$-模，在积集合
\begin{align*}
M=M_1\times M_2\times\cdots\times M_n=\{(x_1,x_2,\cdots,x_n)\mid x_i\in M_i,\,i=1,2,\cdots,n\}
\end{align*}
上定义加法与$R$-数乘运算：
\begin{align}
(x_1,x_2,\cdots,x_n)+(y_1,y_2,\cdots,y_n)=(x_1+y_1,x_2+y_2,\cdots,x_n+y_n),\label{eq:5.2.1}
\end{align}
\begin{align}
a(x_1,x_2,\cdots,x_n)=(ax_1,ax_2,\cdots,ax_n),\quad \forall a\in R.\label{eq:5.2.2}
\end{align}
则$M$构成$R$-模，称为$M_1,M_2,\cdots,M_n$的\textbf{外直和}，记作
\begin{align*}
M=M_1\oplus M_2\oplus\cdots\oplus M_n=\bigoplus\limits_{i=1}^n M_i.
\end{align*}
记$M_i'=\{(0,\cdots,x_i,\cdots,0)\mid x_i\in M_i\}\subseteq M$，则$M_i'$为$M$的子模，且$M_i\cong M_i'$。
\end{definition}

\begin{theorem}
设$M_1,M_2,\cdots,M_n,N$均为$R$-模，$\varphi_i\colon M_i\to N$为模同态，$1\leqslant i\leqslant n$，则存在唯一模同态$\varphi\colon \bigoplus\limits_{i=1}^n M_i\to N$，使得$\varphi(x_i')=\varphi_i(x_i)$，其中$x_i'\in M_i'$，$1\leqslant i\leqslant n$。
\end{theorem}

\begin{proof}

定义映射$\varphi\colon \bigoplus\limits_{i=1}^n M_i\to N$：
\begin{align*}
\varphi(x_1,x_2,\cdots,x_n)=\sum\limits_{i=1}^n \varphi_i(x_i).
\end{align*}
验证模同态条件：
\begin{align*}
\begin{split}
\varphi\big((x_1,\cdots,x_n)+(y_1,\cdots,y_n)\big)
&=\varphi(x_1+y_1,\cdots,x_n+y_n)\\
&=\sum\limits_{i=1}^n \varphi_i(x_i+y_i)=\sum\limits_{i=1}^n \big(\varphi_i(x_i)+\varphi_i(y_i)\big)\\
&=\sum\limits_{i=1}^n \varphi_i(x_i)+\sum\limits_{i=1}^n \varphi_i(y_i)\\
&=\varphi(x_1,\cdots,x_n)+\varphi(y_1,\cdots,y_n),
\end{split}
\end{align*}
对任意$a\in R$，
\begin{align*}
\begin{split}
\varphi\big(a(x_1,\cdots,x_n)\big)&=\varphi(ax_1,\cdots,ax_n)\\
&=\sum\limits_{i=1}^n \varphi_i(ax_i)=\sum\limits_{i=1}^n a\varphi_i(x_i)\\
&=a\sum\limits_{i=1}^n \varphi_i(x_i)=a\varphi(x_1,\cdots,x_n).
\end{split}
\end{align*}
故$\varphi$为模同态，且满足$\varphi(x_i')=\varphi_i(x_i)$，存在性成立。

设$\psi$为另一满足条件的模同态，则对任意$(x_1,\cdots,x_n)=\sum\limits_{i=1}^n x_i'\in \bigoplus\limits_{i=1}^n M_i$，有
\begin{align*}
\psi(x_1,\cdots,x_n)=\psi\left(\sum\limits_{i=1}^n x_i'\right)=\sum\limits_{i=1}^n \varphi_i(x_i)=\varphi(x_1,\cdots,x_n),
\end{align*}
故$\varphi=\psi$，唯一性成立.

\end{proof}

\begin{theorem}
设$M_1,M_2,\cdots,M_n$为$R$-模$N$的子模，若满足
\begin{align}
N=M_1+M_2+\cdots+M_n,\label{eq:5.2.3}
\end{align}
\begin{align}
M_i\cap\left(\sum\limits_{j\neq i} M_j\right)=\{0\},\quad 1\leqslant i\leqslant n,\label{eq:5.2.4}
\end{align}
则映射
\begin{align*}
\varphi\colon \bigoplus\limits_{i=1}^n M_i\to N,\quad \varphi(x_1,x_2,\cdots,x_n)=\sum\limits_{i=1}^n x_i,\ x_i\in M_i
\end{align*}
为模同构。
\end{theorem}

\begin{proof}

记$\varphi_i\colon M_i\to N$为嵌入同态，即$\varphi_i(x_i)=x_i$。
由上一定理，存在模同态$\varphi\colon \bigoplus\limits_{i=1}^n M_i\to N$，使得$\varphi(x_1,\cdots,x_n)=\sum\limits_{i=1}^n x_i$。

由\eqref{eq:5.2.3}，$\varphi$为满同态。
若$\varphi(x_1,\cdots,x_n)=\sum\limits_{i=1}^n x_i=0$，则对任意$1\leqslant i\leqslant n$，有
\begin{align*}
x_i=-\sum\limits_{j\neq i}x_j\in M_i\cap\left(\sum\limits_{j\neq i}M_j\right)=\{0\},
\end{align*}
故$x_i=0$，即$\varphi$为单同态。
因此$\varphi$为模同构.

\end{proof}

\begin{definition}
若$R$-模$N$的子模$M_1,M_2,\cdots,M_n$满足条件\eqref{eq:5.2.3}与\eqref{eq:5.2.4}，则称$N$为$M_1,M_2,\cdots,M_n$的\textbf{内直和}，记作
\begin{align*}
N=M_1\oplus M_2\oplus\cdots\oplus M_n.
\end{align*}
外直和可视为对应内直和，即将$M_i$与$M_i'$等同。
\end{definition}

\begin{proposition}
设$M=\bigoplus\limits_{i=1}^n M_i$，若对每个$i$，$M_i=\bigoplus\limits_{k=1}^{s_i} N_{i_k}$，则
\begin{align*}
M=\bigoplus\limits_{i=1}^n\bigoplus\limits_{k=1}^{s_i} N_{i_k}.
\end{align*}
\end{proposition}

\begin{proposition}
设$M=\bigoplus\limits_{i=1}^n M_i$，取下标划分$1\leqslant r_1<r_2<\cdots<r_s=n$，记
\begin{align*}
N_k=\bigoplus\limits_{t=r_{k-1}+1}^{r_k} M_t,\quad k=1,2,\cdots,s,
\end{align*}
则$M=\bigoplus\limits_{k=1}^s N_k$。
\end{proposition}

\begin{definition}
$R$-模$N$的子模$M_1,M_2,\cdots,M_n$称为\textbf{无关子模}，若满足条件\eqref{eq:5.2.4}。
\end{definition}

\begin{remark}
若$x_1,x_2,\cdots,x_n\in N$为线性无关元素组，则$Rx_1,Rx_2,\cdots,Rx_n$为$N$的无关子模，反之不成立。
\end{remark}

\begin{example}
设$R=\mathbb{F}$为域，$V$为$\mathbb{F}$上线性空间，则元素$x_1,x_2,\cdots,x_n\in V$线性无关当且仅当子空间$\mathbb{F}x_1,\mathbb{F}x_2,\cdots,\mathbb{F}x_n$为无关子模。
\end{example}

\begin{example}
$M=\mathbb{Z}_6$为$\mathbb{Z}$-模，易证$\mathbb{Z}\overline{2}\cap\mathbb{Z}\overline{3}=\{\overline{0}\}$，即$\mathbb{Z}\overline{2},\mathbb{Z}\overline{3}$为无关子模；但$3\overline{2}+2\overline{3}=\overline{0}$，故$\overline{2},\overline{3}$不是线性无关元素。
\end{example}



\end{document}